\chapter{The software environments}
\label{chap:D-environnements}

This appendix gives the versions of the software that produced the results, the platform and the
toolchain that builds the documents. The files in \file{envs/} are authoritative. \Cref{tab:env}
extracts the key packages from them.

\section{The platform}
\label{sec:d-plateforme}

Every run and every published analysis comes from one Mac with an Apple M2 Max chip, under
macOS 27.2, for the arm64 architecture, with Python 3.11.18. The simulations ran on the central
processor, one thread per run. The GRPO runs ran on the chip's graphics processor (the MPS backend of
PyTorch) or on the central processor, depending on the run (\cref{chap:07-modeles-de-langue}). The
original environments come from conda (\texttt{conda create -p sim/env python=3.11}). pip then
installed each package. An ordinary \texttt{venv} environment receives the same packages.

\section{The four environments}
\label{sec:d-envs}

Each file in \file{envs/} is the \texttt{pip freeze} of an environment, taken on 15 November 2026, with the Python
version and the platform at the top. The \file{*.lock} files give the same lists with the SHA-256
fingerprint of each package, for an installation with \texttt{pip install -{}-require-hashes}. Four
\file{*.txt} files, one per environment, repeat the lists that the protocols cite,
\file{sim/requirements.lock} and its siblings, gathered in one folder. The fifth, \file{test.txt},
adds the test packages to the analysis environment.

\begin{table}[htbp]
\centering
\small
\caption{The frozen environments of the project (\file{envs/README.md}).}
\label{tab:d-envs}
\begin{tabularx}{\linewidth}{@{}llX@{}}
\toprule
File & Environment & What it runs \\
\midrule
\file{train.txt} & \file{sim/env} & the training runs (\file{sim/run.py}) and the re-evaluations
  (\file{sim/dose.py}) \\
\file{analysis.txt} & \file{sim/env-an} & every analysis script, \file{docs/check\_numbers.py},
  continuous integration \\
\file{zoo.txt} & \file{sim/env-zoo} & the official RL Baselines3 Zoo pipeline of experiment E6 \\
\file{llm.txt} & \file{sim/env-llm} & GRPO (\file{sim/llm/grpo.py}) and the figures of \file{docs/} \\
\file{test.txt} & on top of \file{analysis.txt} & the unit tests (\file{tests/}) \\
\bottomrule
\end{tabularx}
\end{table}

\gentable{env}

\section{Installation pitfalls}
\label{sec:d-installer}

\begin{itemize}
  \item \textbf{Recreating an environment.} The paths matter, since the scripts and the launchers
        call \file{sim/env/bin/python} and its siblings:
\begin{lstlisting}[language=bash]
python3.11 -m venv sim/env-an
sim/env-an/bin/pip install -r envs/analysis.txt -r envs/test.txt
make verify
\end{lstlisting}
        The other three follow the same pattern. \file{sim/env-zoo} is a \texttt{venv} built on the
        Python of \file{sim/env}.
  \item \textbf{rliable 1.2.0} is distributed only as a source archive. pip builds it in a few
        seconds, since it is pure Python.
  \item \textbf{On Linux}, \file{analysis.txt} installs as is. Continuous integration does so on
        Ubuntu. The last digits of a bootstrap move with the processor: the byte-for-byte comparison
        is required on macOS arm64, the platform of the results. Linux is compared to within
        $10^{-12}$ in relative terms (\cref{sec:ci}).
  \item \textbf{Training on Linux} works with \file{train.txt}, but the runs in the repository come
        from macOS arm64. A new run elsewhere gives other curves for the same random seed: MuJoCo,
        JAX and PyTorch round differently from one processor to another (\cref{sec:reentrainer}).
  \item \textbf{GRPO} downloads \texttt{Qwen/Qwen2.5-0.5B-Instruct} and
        \texttt{unsloth/Llama-3.2-1B-Instruct} from Hugging Face. Since deviation 15 (c) of the GRPO
        strand, \file{grpo.py} pins the exact revision of each model, the one used in our runs.
        \file{docs/figures/build.py} downloads Qwen's tokeniser. Qwen's file
        \file{generation\_config.json} sets a repetition penalty of 1.1, which the published runs
        underwent (deviation 18 of the GRPO strand). The option \texttt{-{}-c7} of \file{grpo.py}
        brings it back to 1.0 for the rerun of deviation 20. Without the option, the code repeats
        the earlier runs identically.
\end{itemize}

\section{The document toolchain}
\label{sec:d-chaine}

Outside the Python environments, building the PDFs and the figures uses the tools of
\cref{tab:d-chaine}. These tools draw and typeset without computing any number. \texttt{make check}
and \texttt{make verify} need only \file{analysis.txt}.

\begin{table}[htbp]
\centering
\small
\caption{The tools that build the documents.}
\label{tab:d-chaine}
\begin{tabularx}{\linewidth}{@{}llX@{}}
\toprule
Tool & Version & Used by \\
\midrule
pandoc & 3.6.1 & \file{docs/build.sh} \\
md2star & 3.13.5 (PyPI) & \file{docs/build.sh} (DOCX and a first PDF) \\
TeX Live 2026: LuaHBTeX, latexmk & 1.24.0, 4.88 & \file{docs/build.sh}, \file{report/build.sh} \\
sprezzature-figures & 2.14.1 (2.14.0 on PyPI) & \file{docs/figures/build.py},
  \file{article/figures/build.py} \\
rsvg-convert (librsvg) & 2.61.3 & conversion of SVG files to PDF and PNG \\
ffmpeg & 8.1.2 & \file{report/tools/new\_figures.py} (frames taken from the E6 videos) \\
git & 2.50.0 & \texttt{make verify} \\
\bottomrule
\end{tabularx}
\end{table}

Two tools outside Python served the analyses of 21 and 22 November 2026. The \texttt{gh} client
2.95.0 read the creation date of the GitHub repositories cited for Q5 (survey deviation 31), one
request at a time. Ollama 0.40.16 serves \texttt{gemma3:12b} locally, the reader of survey deviation
33. Each reading file of this reader records the digest of the model.

This report is built by \texttt{make report}, which runs \file{report/build.sh}: the tables from the
JSON files (\file{report/tools/tables.py}), the figures (\file{report/tools/figures.py}), then two
passes of LuaLaTeX with biber, in French (\file{rapport.tex}) and in English (\file{report-en.tex}).
